\section{Planar equivalent layers located above the sources}
\label{sec:eqlayers_above_sources}


Most of what is known or assumed to be true without any proof about the
equivalent-layer technique can be deduced from the solution of Neumann's 
problem on a plane (equation \ref{eq:Theta-Neumann}).
Note that, for the case in which the observation points 
$\mathbf{r} = (x, y, z)^{\top}$ satisfy $z < z''_{0}$, where $z''_{0}$ defines a 
horizontal plane located above the 3D sources, $\Theta(\mathbf{r})$ (equation 
\ref{eq:scalar-potential-Theta}) has continuous partial derivatives of all orders. 
These derivatives, in turn, can be obtained from equation \ref{eq:Theta-Neumann} 
by differentiating under the sign of integration.
Consider, for example, the partial derivatives of first and second order:
\begin{equation}
	\partial_{\alpha} \Theta(\mathbf{r}) =  
	\int\limits_{-\infty}^{\infty} \int\limits_{-\infty}^{\infty}
	\partial_{z} \Theta(\mathbf{r}'') \:\: 
	\partial_{\alpha} \Psi(\mathbf{r}, \mathbf{r}'') \: dx''dy'' \: , 
	\quad z < z''_{0} \: ,
	\quad \alpha = x, y, z \: ,
	\label{eq:1st-deriv-Theta-Neumann}
\end{equation}
and
\begin{equation}
	\partial_{\alpha\beta} \Theta(\mathbf{r}) =  
	\int\limits_{-\infty}^{\infty} \int\limits_{-\infty}^{\infty}
	\partial_{z} \Theta(\mathbf{r}'') \:\: 
	\partial_{\alpha\beta} \Psi(\mathbf{r}, \mathbf{r}'') \: dx''dy'' \: , 
	\quad z < z''_{0} \: ,
	\quad \alpha, \beta = x, y, z \: ,
	\label{eq:2st-deriv-Theta-Neumann}
\end{equation}
where $\partial_{\alpha} \Theta(\mathbf{r}) \equiv \tfrac{\partial
\Theta(\mathbf{r})}{\partial \alpha}$ and 
$\partial_{\alpha\beta} \Theta(\mathbf{r}) \equiv \tfrac{\partial^{2}
\Theta(\mathbf{r})}{\partial \alpha \, \partial \beta}$.
Note that the derivatives are computed with respect to the coordinates
of $\mathbf{r}$ and, consequently, act only on the kernel of the integral 
defined by equation \ref{eq:Theta-Neumann}.
The following subsections use equations \ref{eq:1st-deriv-Theta-Neumann} and
\ref{eq:2st-deriv-Theta-Neumann} to deduce some closed-form solutions $p(\mathbf{r}'')$ of 
the planar equivalent-layer integral (equation \ref{eq:integral-equation-eqlayer}) 
applied to different potential-field data.

\subsection{Closed-form solutions for gravity data}
\label{subsec:solutions_gravity_data}

Consider the case in which we apply the planar equivalent-layer technique to 
gravity disturbance data \cite[e.g.,][ p. 93]{hofmann-wellenhof-moritz2005}. 
In this case, the potential-field
data are considered as a harmonic function $d(\mathbf{r})$ representing the vertical 
component of the gravitational attraction produced by the 3D sources at the observation
points $\mathbf{r}$. It means that $d(\mathbf{r}) = \partial_{z} \Theta(\mathbf{r})$
(equation \ref{eq:1st-deriv-Theta-Neumann}) and 
$a(\mathbf{r}, \mathbf{r}'') = \partial_{z} \Psi(\mathbf{r}, \mathbf{r}'')$, so that 
the planar equivalent-layer integral (equation \ref{eq:integral-equation-eqlayer}) has
a closed-form solution $p(\mathbf{r}'') = \partial_{z} \Theta(\mathbf{r}'')$.
This result shows that, under the conditions established for the solution of 
Neumann's problem (equation \ref{eq:Theta-Neumann}), there exists a continuous planar
equivalent layer of monopoles that exactly reproduces the observed gravity disturbance
data above the sources.
\citet{dampney1969} was the pioneer in using this result to interpolate gravity data and
a considerable amount of literature has been published on this topic since then.
This solution $p(\mathbf{r}'')$ can also be directly deduced from the
\textit{upward-continuation integral}
\citep[e.g., ][]{peters1949, roy1962, bhattacharyya1967, henderson1970, gunn1975}, 
which represents the solution of Dirichlet’s problem on a plane
(equation \ref{eq:Theta-Dirichlet}).

Consider now the case in which the potential-field data are gravity-gradient tensor 
components. These data are usually described as the partial derivatives of second order 
of the gravitational potential. Hence, by using the second derivatives of 
$\Theta(\mathbf{r})$ (equation \ref{eq:2st-deriv-Theta-Neumann}), we obtain for the 
$\alpha\beta$-component of gravity tensor that
$d(\mathbf{r}) = \partial_{\alpha\beta} \Theta(\mathbf{r})$ and 
$a(\mathbf{r}, \mathbf{r}'') = \partial_{\alpha\beta} \Psi(\mathbf{r}, \mathbf{r}'')$, 
so that the planar equivalent-layer integral 
(equation \ref{eq:integral-equation-eqlayer}) has a closed-form solution 
$p(\mathbf{r}'') = \partial_{z} \Theta(\mathbf{r}'')$.
This result shows that the same continuous planar equivalent layer of monopoles that 
exactly reproduces the observed gravity disturbance data also reproduces the 
gravity gradient tensor data above the sources.
This result has been used by the scientific community for processing gravity-gradient
data \citep{barnes-lumley_2011} or converting gravity disturbance into the
gravity-gradient tensor \citep{piauilino_etal2019_seg, li_etal2021}, for example.
However, researchers have not treated this topic in much detail and it is usually assumed
to be true without any proof.

\subsection{Closed-form solutions for magnetic data}
\label{subsec:solutions_magnetic_data}